\subsection{逐点排除法}
\label{subsec:lab4-point-exclusion}

\paragraph{基本原理}

简单来说就是正难则反，我们如果能求出点集的补集，那也能刻画出点集

也就是说我们排除掉不可达的点，剩下的就是可达的点了，那怎么排除呢？

我们知道两个点重合就意味他们的距离为0，如果距离大于0，那就意味着两个点不是同一个点，可以用这个特性来排除

那距离首先想到的就是Euclid距离

设不可达点集为
\[
Q=\{0,1\}^{8}\setminus V.
\]
我们
对于任意$q=(q_1,\ldots,q_8)\in Q$
\[
d(x,q)
= \sqrt{\sum_{i=1}^8 (x_i-q_i)^2}
\]
我们知道二进制不管怎么加减无非就是0,1,-1,那么也就是说\((x_i-q_i)^2 = |x_i-q_i|\)

那我们再按$q_i$的值来分的话

如果$q_i=0$,那$|x_i-q_i| = |x_i| = x_i$,

如果$q_i=1$,那$|x_i-q_i| = |x_i - 1| = 1-x_i$,因为$x_i-1 \leq 0$

那么也就是说
\[
d(x,q)
=\sqrt{\sum_{j:q_j=0}x_j+\sum_{j:q_j=1}(1-x_j)}
\]

那我们希望是整数，干脆就不要开根号了，反正没什么区别

那这其实就是Hamming距离

对于任意$q=(q_1,\ldots,q_8)\in Q$，定义$x$与$q$的Hamming距离
\[
d(x,q)
=\sum_{j:q_j=0}x_j+\sum_{j:q_j=1}(1-x_j).
\]

\paragraph{数学建模}
\[
\begin{aligned}
\min\quad &0\\
\text{s.t.}\quad
&\sum_{j:q_j=0}x_j+\sum_{j:q_j=1}(1-x_j)\ge1,
&&q\in Q,\\
&x_j\in\{0,1\},
&&j=1,\ldots,8.
\end{aligned}
\]
也可以进行移项
\[
\sum_{j:q_j=0}x_j-\sum_{j:q_j=1}x_j
\ge1-\sum_{j=1}^{8}q_j,
\qquad q\in Q
\]

例如，输入为$0$、输出为$1$对应的点为$q=0000\,0001$,我们如果要排除它，就可以写成
\[
x_1+x_2+x_3+x_4+x_5+x_6+x_7+(1-x_8)\ge1.
\]


\paragraph{正确性分析}

对于任意不可达点 $q\in Q$，若 $x=q$，则
\[
d(x,q)=\sum_{j:q_j=0}x_j+\sum_{j:q_j=1}(1-x_j)=0,
\]
不满足约束 $d(x,q)\ge1$，所以不可达点都被排除了

反之，若 $x\ne q$，由于 $x,q\in\{0,1\}^8$，
则至少存在一个坐标 $j$ 使得 $x_j\ne q_j$，
因此 $d(x,q)\ge1$，也就是说只排除$q$这一个点

对于任意可达点 $x\in V$，由于 $V$ 与 $Q$ 不相交，
因此对于所有 $q\in Q$，均有 $x\ne q$，
从而满足全部约束，不会有可达点遗漏

反过来，若 $x\notin V$，则根据
$Q=\{0,1\}^8\setminus V$，必有 $x\in Q$。
此时模型中存在针对 $q=x$ 的约束，而
$d(x,q)=0$，与 $d(x,q)\ge1$ 矛盾,所以不会引入任何不可达点。

综上，模型的二进制可行解集合恰好为 $V$，即
\[
\left\{
x\in\{0,1\}^8\;\middle|\;
\sum_{j:q_j=0}x_j+\sum_{j:q_j=1}(1-x_j)\ge1,
\ \forall q\in Q
\right\}=V.
\]

\paragraph{用于差分与线性传播}

差分建模时，将DDT零位置对应的八位向量作为$Q$；线性建模时，将LAT零位置对应的八位向量作为$Q$。两种情况下均只使用原始变量$x$，不需要额外的选择变量。
